% ============================================================
%  Chapter 6 — How Messages and Videos Travel
%  The Secret Life of the Internet
% ============================================================

\chapter{How Messages and Videos Travel}

\chapteropening{%
  A video is enormous.  How can something so big travel so fast?
  The secret is that it does not travel whole —
  it is broken into thousands of tiny pieces called \textbf{packets}!
}
\chapterbadge{Packets on the Move}

% --- Opening story ---
\section*{Dot's Secret: There Is More Than One Dot}

Dot had been humble about something.

\begin{center}
  \Large\bfseries\color{InternetBlue}
  ``I am not alone,'' Dot admitted.
\end{center}

When Sam sent a photo to a friend, it did not travel as one big lump.
It was cut into hundreds of small pieces.
Each piece was wrapped in an envelope with an address written on it.
And each envelope was a \emph{packet}.

Suddenly the air around Dot filled with hundreds of identical little glowing dots,
all labelled with numbers: Packet 1, Packet 2, Packet 3\ldots

\character{Dot}{We are all parts of the same photo!
  We might travel different routes to get there,
  but we will all arrive and get put back together at the end.
  It is like posting a jigsaw puzzle one piece at a time.}

\sectionrule

% --- What is a packet? ---
\section*{What Is a Packet?}

A \term{packet} is a small chunk of data.

\begin{funfact}
A single 4K video stream sends about 25 Megabits of data every second — that's roughly 3,000 packets flying through the internet every second just for your film!
\end{funfact}

When you send anything over the internet — a message, a photo, a video — it is split into packets.
Each packet contains:
\begin{itemize}
  \item A piece of the data (part of the photo, video, or message).
  \item A \term{header} — information about where the packet came from, where it is going,
        and what number it is in the sequence.
\end{itemize}

\begin{picturethis}
  Imagine sending a 1000-piece jigsaw puzzle through the post.
  You cannot fit all 1000 pieces in one envelope.
  So you put 10 pieces in each envelope, label each envelope with a number,
  and post them all.
  At the other end, the person sorts the envelopes by number
  and puts the puzzle back together.
  \textbf{That is exactly how packets work.}
\end{picturethis}

\sectionrule

% --- Routing ---
\section*{How Do Packets Find Their Way?}

Packets do not all have to take the same route.
Along the internet highway, special computers called \term{routers} read the address on each packet
and direct it toward the destination.

If one route is busy or broken, a router sends the packet a different way.

\begin{quickmap}
  \textbf{The packet journey:}
  \begin{enumerate}[label=\textbf{\arabic*.}]
    \item A photo is split into 200 packets, each numbered 1 to 200.
    \item All 200 packets set off from Sam's device.
    \item Routers along the way read each packet's address and forward it.
    \item Some packets take a northern route; others take a southern route.
    \item All 200 packets arrive at the destination device.
    \item The device sorts them back into order using the numbers.
    \item The device reassembles the original photo. Done!
  \end{enumerate}
\end{quickmap}

\begin{diagram}[Packets split apart, travel along different routes, and reassemble.]
  \resizebox{\linewidth}{!}{%
  \begin{tikzpicture}[>=Stealth]
    \node[draw=InternetBlue, rounded corners, fill=SkyBlue!25, minimum width=2.9cm, minimum height=1.3cm] (original) at (-4.8,0) {Original Photo};

    \node[draw=WarmOrange!80!black, rounded corners, fill=WarmOrange!20, minimum width=1.9cm, minimum height=0.7cm] (p1) at (-0.7,1.5) {Packet 1};
    \node[draw=LeafGreen!80!black, rounded corners, fill=LeafGreen!20, minimum width=1.9cm, minimum height=0.7cm] (p2) at (-0.7,0.5) {Packet 2};
    \node[draw=SoftPurple!80!black, rounded corners, fill=SoftPurple!20, minimum width=1.9cm, minimum height=0.7cm] (p3) at (-0.7,-0.5) {Packet 3};
    \node[draw=CoralRed!80!black, rounded corners, fill=CoralRed!16, minimum width=1.9cm, minimum height=0.7cm] (p4) at (-0.7,-1.5) {Packet 4};

    \node[draw=InternetBlue!80!black, rounded corners, fill=SunYellow!25, minimum width=3.2cm, minimum height=1.3cm] (reassembled) at (4.4,0) {Reassembled Photo};

    \draw[thick,->,InternetBlue] (original) -- (p1);
    \draw[thick,->,InternetBlue] (original) -- (p2);
    \draw[thick,->,InternetBlue] (original) -- (p3);
    \draw[thick,->,InternetBlue] (original) -- (p4);

    \draw[dotted,thick,->,InternetBlue] (p1) .. controls (1.4,2.2) and (2.3,1.0) .. (reassembled);
    \draw[dotted,thick,->,InternetBlue] (p2) .. controls (1.2,1.0) and (2.3,0.5) .. (reassembled);
    \draw[dotted,thick,->,InternetBlue] (p3) .. controls (1.5,-0.6) and (2.3,-0.2) .. (reassembled);
    \draw[dotted,thick,->,InternetBlue] (p4) .. controls (1.4,-1.8) and (2.3,-1.0) .. (reassembled);

    \node[align=center,font=\small\bfseries\color{WarmOrange}] at (-2.6,2.35) {Split into packets};
    \node[align=center,font=\small\bfseries\color{SoftPurple}] at (2.0,2.35) {Different routes};
  \end{tikzpicture}%
  }
\end{diagram}

\character{Dot}{Even if I arrive before Packet 50, the device waits for everyone.
  If a packet gets lost, the device asks for it to be sent again.
  Nobody gets left behind!}

\sectionrule

% --- TCP/IP ---
\section*{The Rules of the Road: TCP/IP}

The internet follows strict rules so that packets behave properly.
The most important set of rules is called \term{TCP/IP}.

\begin{itemize}
  \item \textbf{IP (Internet Protocol)} — gives every device on the internet its own address
        (called an \term{IP address}) and makes sure packets have the right destination written on them.
  \item \textbf{TCP (Transmission Control Protocol)} — makes sure all the packets arrive safely.
        If a packet is missing, TCP asks for it to be sent again.
\end{itemize}

\begin{picturethis}
  Think of IP as the postal address system — every device has an address.
  Think of TCP as a very careful delivery service that checks every item off a list
  and will not sign off until every single package has arrived safely.
\end{picturethis}

\sectionrule

% --- Video streaming ---
\section*{What About Videos?}

A video is made of thousands of still images flashed one after another,
plus an audio track.
A one-minute video can be made up of millions of packets.

When you \term{stream} a video, the packets are not all downloaded first.
Instead, packets keep arriving just fast enough to play the video smoothly.
That is why if your Wi-Fi is slow, the video might pause —
it is waiting for the next wave of packets to arrive!

\sectionrule

\begin{newwords}
  \begin{description}[labelwidth=3.2cm, leftmargin=3.4cm]
    \item[\term{Packet}]   A small chunk of data sent across the internet, with a numbered label and address.
    \item[\term{Header}]   The label on a packet that says where it came from, where it is going, and its sequence number.
    \item[\term{Router}]   A device that reads packets and directs them toward their destination.
    \item[\term{IP}]       Internet Protocol — the addressing system of the internet.
    \item[\term{TCP}]      Transmission Control Protocol — the system that makes sure all packets arrive correctly.
    \item[\term{Streaming}] Receiving and playing data (like a video) as it arrives, rather than waiting for it all first.
  \end{description}
\end{newwords}

\sectionrule

\begin{tryit}
  \textbf{Be the packet sorter!}

  Below are 8 packets out of order.
  Number them correctly, then ``play'' the message:

  \begin{center}
    \begin{tabular}{|c|l|}
      \hline
      \textbf{Packet \#} & \textbf{Fragment} \\
      \hline
      ?  & ``the internet!'' \\
      ?  & ``Hello, my name is'' \\
      ?  & ``I love learning'' \\
      ?  & ``Dot, and'' \\
      ?  & ``how'' \\
      ?  & ``about'' \\
      ?  & ``things'' \\
      ?  & ``like'' \\
      \hline
    \end{tabular}
  \end{center}

  \emph{Correct order: ``Hello, my name is Dot, and I love learning about things like the internet!''}

  Imagine doing that with millions of pieces of a video!
\end{tryit}

\sectionrule

\begin{bigidea}
  Information on the internet travels as tiny pieces called \textbf{packets}.
  Each packet finds its own route, and they are all put back together
  in the right order when they arrive.
\end{bigidea}

\character{Dot}{Amazing!  Now you know my biggest secret — I'm made of packets!
  Next up, let's meet Scout Search and find out how search engines
  help you find exactly what you are looking for in the giant internet library.}
